% ---
% id: protocol-compiler/req:operator-form
% kind: req
% title: "Operators on the measured system held as block lists"
% summary: >-
%   Every operator on the measured system is held and emitted as its block
%   list. Where there are several sectors its direct sum is formed only as the
%   argument of chain condition (C4), decided there and discarded.
% prominence: standard
% status: accepted
% uses: [def:declaration, def:chaincond]
% assumes: []
% refs: []
% provenance: newly-derived
% lean:
% ---

By Definition~\ref{def:declaration} an operator on the \tref{measured system} is its
\tref{block list}, so sector block-diagonality holds by type and the declaration
holds no array larger than one sector block.

\begin{requirement}\label{req:operator-form}
The engine holds and emits every operator on $\Qsys$ as its
\tref{block list} against the sector map $\Smap$
(Definition~\ref{def:declaration}). Where $|\SecSet|>1$ it does not
materialize the direct sum $\bigoplus_sF^{(s)}$ of such a list, with one
exception: condition (C4) of Definition~\ref{def:chaincond} forms it as the
argument of that one decision, decides it and discards it. The convention
\iface{operator_form} states this rule and its scope.
\end{requirement}
